%template para tesis del Instituto Tecnológico Superior de Misantla
%Creado por: José Arcángel Salazar Delgado
\documentclass[dvips,12pt]{report}
\usepackage[utf8]{inputenc}
\usepackage[spanish]{babel}
\usepackage[pdftex]{graphicx}
\usepackage{url}
\usepackage{enumerate} 
\usepackage{cite}
\usepackage{apacite}
\usepackage{datetime2}
\usepackage{fancyhdr}
\usepackage{placeins}
\usepackage{float}
\usepackage{booktabs}
\usepackage{listings}


\graphicspath{ {images/} }

\setlength{\oddsidemargin}{0.25in}
\setlength{\textwidth}{6.5in}
\setlength{\topmargin}{0in}
\setlength{\textheight}{8.5in}

\begin{document}

\thispagestyle{empty}

\begin{minipage}[c][0.1\textheight][c]{0.25\textwidth}
\begin{center}
    \includegraphics[width=4cm, height=4cm]{logo-misantla.jpg}
\end{center}
\end{minipage}
\begin{minipage}[c][0.05\textheight][t]{0.7\textwidth}
\begin{center}
    {\scshape Instituto Tecnol\'ogico Superior de Misantla}
    \vspace{.3cm}
    \hrule height2.5pt
    \vspace{.1cm}
    \hrule height1pt
    \vspace{.3cm}
    {\scshape  Maestr\'ia en sistemas computacionales}
\end{center}
\end{minipage}

\begin{minipage}[c][0.5\textheight][t]{0.25\textwidth}
	\begin{center}
	\hskip2pt

	\vrule width1pt height18cm 
	\hskip1mm
	\vrule width2.5pt height18cm
    \hskip1mm
    \vrule width1pt height18cm \\
    \end{center}
\end{minipage}
\begin{minipage}[c][0.6\textheight][t]{0.65\textwidth}
  \begin{center}
    \vspace{2cm}
    {\huge \scshape {T\'itulo de la tesis}}
    \vspace{2cm}

    \makebox[5cm][c]{\LARGE TESIS}  \\[8pt]
    QUE PARA OBTENER EL T\'iTULO DE:\\[5pt]
    {\large \textbf{{t\'itulo}}}\\[40pt]            
    PRESENTA:\\[5pt]
    \textbf{{Nombre del tesista}}

    \vspace{1cm}

    {\small DIRECTOR DEL TRABAJO:\\ {Nombre del director del trabajo}}

    \vspace{0.5cm}

    {Lugar,}{ }{Fecha}
  \end{center}
\end{minipage}

%template para tesis del Instituto Tecnológico Superior de Misantla
%Creado por: José Arcángel Salazar Delgado
\addcontentsline{toc}{chapter}{Agradecimientos}
\thispagestyle{plain}
\begin{center}
	\Large    
	\vspace{0.9cm}
	\textbf{Agradecimientos}
\end{center}
%Abstract en ingles
Lorem ipsum dolor...
%Plantilla para tesis del Instituto Tecnológico Superior de Misantla
%Creada por: José Arcángel Salazar Delgado
%Supervisó: Dr. Eddy Sánchez de la Cruz
%Vo.Bo.: MSC. Galdino Martínez Flores.
%Versión: 1.0
\addcontentsline{toc}{chapter}{Abstract}
\thispagestyle{plain}
\begin{center}
    \Large    
    \vspace{0.9cm}
    \textbf{Abstract}
\end{center}
%Abstract en ingles
En esta sección se resume en inglés el contenido del trabajo de investigación.
%Plantilla para tesis del Instituto Tecnológico Superior de Misantla
%Creada por: José Arcángel Salazar Delgado
%Supervisó: Dr. Eddy Sánchez de la Cruz
%Vo.Bo.: MSC. Galdino Martínez Flores.
%Versión: 1.0
\addcontentsline{toc}{chapter}{Resumen}
\thispagestyle{plain}
\begin{center}
	\Large    
	\vspace{0.9cm}
	\textbf{Resumen}
\end{center}
%Abstract en ingles
En esta sección se resume en español el contenido del trabajo de investigación.

\tableofcontents % indice de contenidos
\listoffigures %indice de figuras
\renewcommand\listtablename{Índice de tablas}
\listoftables %indice de tablas

\renewcommand\lstlistingname{Índice de código}
\renewcommand\lstlistlistingname{Índice de código}
\lstlistoflistings

\bibliographystyle{apacite}

%template para tesis del Instituto Tecnológico Superior de Misantla
%Creado por: José Arcángel Salazar Delgado
\chapter{Planteamiento del problema}
\section{Ejemplo de sección}

Praesent in sapien. Lorem ipsum dolor sit amet, consectetuer adipiscing 
elit. Duis fringilla tristique neque. Sed interdum libero ut metus. 
Pellentesque placerat. Nam rutrum augue a leo. Morbi sed elit sit amet 
ante lobortis sollicitudin.

Ejemplo de una lista enumerada:

\begin{enumerate}
	\item uno
	\item dos
	\item tres
\end{enumerate}

Praesent in sapien. Lorem ipsum dolor sit amet, consectetuer adipiscing 
elit. Duis fringilla tristique neque. Sed interdum libero ut metus. 
Pellentesque placerat. Nam rutrum augue a leo. Morbi sed elit sit amet 
ante lobortis sollicitudin.

Ejemplo de una lista con viñetas

\begin{itemize}
 \item prueba 
 \item prueba 
 \item prueba
\end{itemize}

Ejemplo de insertar una figura:
\begin{figure}[H]
	\caption{Logo de la institución.}
	\centering
	\includegraphics[width=4cm, height=4cm]{logo-misantla.jpg}
\end{figure}

Praesent in sapien. Lorem ipsum dolor sit amet, consectetuer adipiscing 
elit. Duis fringilla tristique neque. Sed interdum libero ut metus. 
Pellentesque placerat. Nam rutrum augue a leo. Morbi sed elit sit amet 
ante lobortis sollicitudin.

\section{Otra sección}

Praesent in sapien. Lorem ipsum dolor sit amet, consectetuer adipiscing 
elit. Duis fringilla tristique neque. Sed interdum libero ut metus. 
Pellentesque placerat. Nam rutrum augue a leo. Morbi sed elit sit amet 
ante lobortis sollicitudin.

Ejemplo de una tabla:

\begin{table}[h!]
	\centering
	\caption{Caption for the table.}
	\label{tab:table1}
	\begin{tabular}{ccc}
		\toprule
		Some & actual & content\\
		\midrule
		prettifies & the & content\\
		as & well & as\\
		using & the & booktabs package\\
		\bottomrule
	\end{tabular}
\end{table}

Otro ejemplo de tablas



Ejemplo de citado \cite{Vickrey1961}.

Ejemplo de citado 2 \citeA{Vickrey1961}

Ejemplo de citado 3 \cite{baggio}

Ejemplo de citado 4 \citeA{baggio}




%Plantilla para tesis del Instituto Tecnológico Superior de Misantla
%Creada por: José Arcángel Salazar Delgado
%Supervisó: Dr. Eddy Sánchez de la Cruz
%Vo.Bo.: MSC. Galdino Martínez Flores.
%Versión: 1.0
\chapter{Estado del arte}
%template para tesis del Instituto Tecnológico Superior de Misantla
%Creado por: José Arcángel Salazar Delgado
\chapter{Marco Teórico}
%template para tesis del Instituto Tecnológico Superior de Misantla
%Creado por: José Arcángel Salazar Delgado
%Versión: 1.0
\chapter{Experimentos y resultados}
%Plantilla para tesis del Instituto Tecnológico Superior de Misantla
%Creada por: José Arcángel Salazar Delgado
%Supervisó: Dr. Eddy Sánchez de la Cruz
%Vo.Bo.: MSC. Galdino Martínez Flores.
%Versión: 1.0
\chapter{Análisis de resultados}
%template para tesis del Instituto Tecnológico Superior de Misantla
%Creado por: José Arcángel Salazar Delgado
\chapter{Conclusiones y trabajos futuros}
\bibliography{biblio}{}	
\end{document}